%% ma: L12-B14-0097
%% lop: 12
%% chuong: 5
%% bai: 14
%% dang: 12.14.29
%% loai: DS
%% muc: [NB, TH, TH, TH]
%% dap_an: "ĐSĐS"
%% nguon: "Ảnh giáo viên gửi 07/11/2026 (Tư Duy Mở - tuduymo.com, đề tự luyện Phương trình mặt phẳng, đề lập phương trình mặt phẳng), Câu 13"
%% kien_thuc: "Bài 14 (phương trình mặt phẳng tổng quát, các trường hợp riêng, vị trí của mặt phẳng với trục tọa độ)"
%% trang_thai: chinh-thuc
%% ghi_chu: "Tài liệu gốc không kèm đáp án; đáp án do Claude giải và kiểm bằng sympy. Dạng đã duyệt 07/11/2026."
\begin{ex}
Trong không gian $Oxyz$, xét các mệnh đề về các mặt phẳng đặc biệt sau.
\choiceTF
{\True Phương trình tổng quát của mặt phẳng $(Oyz)$ là $x=0$.}
{Nếu mặt phẳng $(\alpha)$ có phương trình $x+2z-2=0$ thì $(\alpha)\parallel(Ozx)$.}
{\True Nếu mặt phẳng $(\beta)$ có phương trình $3x+4z-1=0$ thì $Oy\parallel(\beta)$.}
{Nếu mặt phẳng $(\gamma)$ có phương trình $2y-5z=0$ thì $Ox\parallel(\gamma)$.}
\loigiai{a) $(Oyz)$ gồm các điểm có hoành độ bằng $0$: $x=0$. Đúng.

b) $(Ozx)\colon y=0$ có vectơ pháp tuyến $(0;1;0)$, không cùng phương với $(1;0;2)$. Sai.

c) $\vec j\cdot(3;0;4)=0$ và $O\notin(\beta)$ (vì $-1\ne0$) nên $Oy\parallel(\beta)$. Đúng.

d) $\vec i\cdot(0;2;-5)=0$ nhưng $O\in(\gamma)$ nên $Ox$ nằm trong $(\gamma)$, không song song. Sai.}
\end{ex}
